% Limitations (0.4 p); full version in Appendix C. Sources: docs/CHANNEL_MODEL_LIMITATIONS.md, EVIDENCE.md
\section{Limitations and Threats to Validity}
\label{sec:limits}

The radio is real 5G software, but propagation is injected: the channel has
large-scale loss and correlated shadowing but no multipath, terminal height is
fixed at 0.5\,m, and absolute availability is internally consistent rather than
externally calibrated. The OAI UE may buffer uplink data differently from a
commercial modem, and constant-rate video is a worst case for the queue;
whether C1 transfers is the first question for hardware. Every
between-terrain claim rests on four terrains. The forecasting and
altitude-effect data predate the cap, so their forest labels are partly queue
backlog, and the dead-process rule removes only about half of the dead cycles
the monitor recorded. Throughput was not measured with active probes during the
campaigns, because they would have congested the links under test; the
emulator's radio thread, near one full core, may cap the cell's total uplink
\todo{measured capacity from the 15 Nov sweep}.
Placement comparisons use one 90-min run per arm with within-run
intervals, and the robots' safety guards react to link state, so the same seed
does not give identical patrols across arms. Appendix~\ref{app:threats} gives
the full list.
\todo{revise the last-but-one sentence once the three-seed repeats are in}
